\section{\texorpdfstring{Proofs for Sections~\ref{sec:general-start}
and~\ref{sec:frontier-modes}}{Proofs for the start and the field}}\label{app:start-field}

\subsection{The start weights}

We prove Lemma~\ref{lem:start-weights}. Throughout $1/d<p<1$, so $\sigma=p-r>0$
and $v=r/\sigma=u/(1-u)$. The refresh form in Proposition~\ref{prop:dpmf}(a) is linear in $\mu$, and with
$\mu=\delta_a$ it gives $p^{N-1}S^{(a)}_n(u)$. So, with sums over
$1\le m_b\le n_b$ for $b\in F$ and $M=\sum_bm_b$,
\begin{equation}\label{eq:startw-weights}
 \omega_a(n)=\E_w\Bigl[\frac{m_a}M\Bigr],\qquad
 w(m)\propto\frac{M!}{\prod_bm_b!}\,v^M\prod_{b\in F}\binom{n_b-1}{m_b-1}.
\end{equation}

\begin{proof}[Proof of Lemma~\ref{lem:start-weights}]
(1) Exchanging the letters $a$ and $b$ is a bijection from the words with counts
$n$ that start with $a$ to those that start with $b$. It keeps the counts and
the number of changes, so $S^{(a)}_n=S^{(b)}_n$. If $k$ divides $N$, a balanced
$n$ has equal coordinates, and the $k$ weights are equal and sum to 1.

(2) We may take $N\ge N_0$ so large that $p>1/d$ and $\sigma\ge\frac12$. Let $n$
be balanced and $\nu=\lfloor N/k\rfloor$, so each $n_b$ is $\nu$ or $\nu+1$. By (1)
the weights take a value $\omega_+$ where $n_b=\nu+1$ and a value $\omega_-$ where
$n_b=\nu$. We show that $0\le\omega_+-\omega_-\le\delta$ with
$\delta=\frac2\nu(\E_wM+\E_wM^2)$. Since the weights average to $1/k$, each is
then within $\delta$ of $1/k$, and
$|\sum_a\mu_a\omega_a(n)-\mu(F)/k|\le\mu(F)\delta$.

Let $n_a=\nu+1$ and $n_b=\nu$, and let $m'$ be $m$ with $m_a$ and $m_b$ exchanged.
All factors of $w(m)$ other than $\binom{\nu}{m_a-1}\binom{\nu-1}{m_b-1}$ are
symmetric in $m_a$ and $m_b$. From $\binom\nu j/\binom{\nu-1}j=\nu/(\nu-j)$ we get
$w(m')=w(m)(\nu-m_a+1)/(\nu-m_b+1)$ when $m_a\le\nu$, and $w(m')=0$ when $m_a=\nu+1$.
In both cases $w(m)-w(m')=w(m)(m_a-m_b)/(\nu-m_b+1)$ when $m_a>m_b$. Pairing $m$
with $m'$ in~\eqref{eq:startw-weights} gives, with $Z=\sum_mw(m)$,
\[
 \omega_a-\omega_b=\frac1Z\sum_{m_a>m_b}\bigl(w(m)-w(m')\bigr)\frac{m_a-m_b}M
 =\frac1Z\sum_{m_a>m_b}w(m)\frac{(m_a-m_b)^2}{M(\nu-m_b+1)}\ge0.
\]
Here $(m_a-m_b)^2/M\le M$. If $m_b\le\nu/2$ the last factor is at most $2/\nu$. If
$m_b>\nu/2$, then $M>\nu/2$ and the term is at most $w(m)M<w(m)\,2M^2/\nu$. This
proves the bound $\delta$.

It remains to bound the moments. Since
$(m-1)!\binom{n_b-1}{m-1}/n_b^{m-1}=\prod_{j<m}(1-j/n_b)_+$, we have
$w(m)\propto c(m)A(m)$ for all $m$ with every $m_b\ge1$, where
\[
 c(m)=M!\prod_b\frac{(n_bv)^{m_b}}{m_b!(m_b-1)!},\qquad
 A(m)=\prod_b\prod_{j<m_b}\Bigl(1-\frac j{n_b}\Bigr)_+\in[0,1],
\]
and Lemma~\ref{lem:clipped-product} gives $1-A(m)\le\sum_bm_b^2/(2n_b)\le M^2/\nu$.
Write $M!=\int_0^\infty e^{-t}t^M\,dt$. Since
$\sum_{m\ge1}y^m/(m!(m-1)!)=y\Phi(y)$, this gives
$\E_c[(M+1)(M+2)]=\langle t^2\rangle$, where $\langle\cdot\rangle$ is the mean
under the density proportional to $e^{-t}t^k\prod_b\Phi(n_bvt)$, as in the
proof of Theorem~\ref{thm:uniform-simplex-majorant}. The integration by parts
there, with $t^{k+2}$ in place of $t^{k+1}$, gives
$\langle t^2\rangle=(k+2)\langle t\rangle+\sum_b\langle t\psi(n_bvt)\rangle$. Put
$c_1=\sqrt v\sum_b\sqrt{n_b}$. That proof gives
$\langle t\rangle\le2(k+1)+c_1^2$, and $\psi(x)\le\sqrt x$ and
$\langle t^{3/2}\rangle\le\langle t^2\rangle^{3/4}$ give
$X=\langle t^2\rangle\le(k+2)(2k+2+c_1^2)+c_1X^{3/4}$. If $c_1X^{3/4}\ge X/2$,
then $X\le16c_1^4$. So $X\le C_2(1+c_1^2)^2$, where $C_2$ depends only on $k$.
Here $c_1^2\le kNv=kT/\sigma\le2kT\le2kCL$, so
$\E_cM^2\le\E_c[(M+1)(M+2)]\le\nu/2$ for $N\ge N_0$. Hence
\[
 \E_wM\le\E_wM^2\le\frac{\E_cM^2}{1-\E_cM^2/\nu}\le2C_2(1+2kT)^2.
\]
With $\nu\ge N/(2k)$ this gives $\delta\le C_1(1+T)^2/N$.
\end{proof}

\subsection{Confinement and the field}

The next lemma is the Gaussian profile of the ratio to a vertex near the center
of a face. We use it for Theorem~\ref{thm:weakfield-phases}.

\begin{lemma}[Gaussian confinement]\label{lem:confinement}
Fix a support $S$ with $|S|=k\ge2$ and a compact set $\mathcal J$ of values of
$t$, and put $T=L-\frac12\log L+t$. For a positive count vector $n$ on $S$ put
$y=\sqrt L\,(n/N-\frac1k\mathbf1)$, a vector in the hyperplane
$\{y\in\mathbb R^S:\sum_iy_i=0\}$.
\begin{enumerate}[(a)]
\item Uniformly for $t\in\mathcal J$ and $y$ in a compact subset of this
hyperplane,
\[
 \log S_n(u)=(k-1)(t-a_k)-\frac{k^2}4|y|^2+o(1).
\]
\item For every $R>0$, uniformly for $t\in\mathcal J$ and all $n$ with $|y|>R$,
\[
 \log S_n(u)\le(k-1)(t-a_k)-\frac{k^2}4R^2+o(1).
\]
\end{enumerate}
\end{lemma}

\begin{proof}
(a) Put $\alpha=n/N=\frac1k\mathbf1+y/\sqrt L$. For $|y|$ bounded,
$\alpha_i\ge1/(2k)$ for large $N$, and Taylor expansion of
$\sqrt{1/k+x}$ in $x$, with $\sum_iy_i=0$, gives
\[
 \Bigl(\sum_i\sqrt{\alpha_i}\Bigr)^2=k-\frac{k^2}{4L}|y|^2+O(L^{-3/2}),
 \qquad D(\alpha)=D_k(1+o(1)).
\]
Insert these into Theorem~\ref{thm:simplex-interior} with $z=Nu=T/p$, where
$z-T=O(L^2/N)$. Then
$(A^2-1)z=(k-1)z-\frac{k^2}4|y|^2+o(1)$ and
$z+\frac12\log z-L=t+\frac12\log(z/L)+o(1)=t+o(1)$. Since
$\log D_k=-(k-1)a_k$, this gives (a).

(b) Starting from $n$, move one unit from a coordinate $b$ to a coordinate $a$
with $n_b\ge n_a+2$, and repeat. The moves keep the support and end at a
balanced vector, and by Theorem~\ref{thm:simplex-balance} each move strictly
increases $S_n(u)$. A move changes $|y|^2$ by $\frac{2L}{N^2}(n_a-n_b+1)<0$ and
changes $y$ by a vector of length $\sqrt{2L}/N$. The balanced end has
$|y|=O(\sqrt L/N)$. So if $|y|>R$ at the start, the first vector on the path
with $|y|\le R$ has $|y|=R+o(1)$, and by (a) its value of $\log S_n$ is
$(k-1)(t-a_k)-k^2R^2/4+o(1)$, uniformly. The starting vector has a smaller value.
\end{proof}

\begin{proof}[Proof of Theorem~\ref{thm:weakfield-phases}]
Fix compact sets of $t$ and $h$, and a support $S$ with $|S|=k\ge2$. Relative to
$V_N$, a bin with support $S$ has weight $e^{h\cdot\alpha}S_n(u)$ with
$\alpha=n/N$. Choose $R$ with $k^2R^2/4>\max_ih_i-\bar h_S+1$ on the compact set.
Since $h\cdot\alpha\le\max_{i\in S}h_i$, Lemma~\ref{lem:confinement}(b) shows
that bins with $|y|>R$ have log weight at most $\ell_S(t,h)-1+o(1)$. On
$|y|\le R$ we have $h\cdot\alpha=\bar h_S+O(R/\sqrt L)$, and
Lemma~\ref{lem:confinement}(a) gives the log weight
$\ell_S(t,h)-\frac{k^2}4|y|^2+o(1)$. A balanced bin has $|y|=O(\sqrt L/N)$. So the
largest weight on $S$ is $e^{\ell_S(t,h)+o(1)}$, while a vertex $i$ has weight
exactly $e^{h_i}$. There are finitely many supports. This
proves~\eqref{eq:weakfield-maximum}, and a positive gap to a lower score excludes
that support for large $N$. If the maximizers on a selected support had
$|y|\ge\eta>0$ along a subsequence, their log weight would be at most
$\ell_S-k^2\eta^2/4+o(1)$, below that of a balanced bin. This proves the
location statement.

Among supports of size $k$ the largest $\bar h_S$ is the average of the $k$
largest fields, which gives~\eqref{eq:weakfield-lines}. If $k_1$ wins at $t_1$
and $k_2$ at $t_2>t_1$, adding the 2 comparisons gives
$(k_1-k_2)(t_2-t_1)\le0$. So the winning size is nondecreasing away from ties.

To realize a given $2\le k<d$, set $h_1=\cdots=h_k=0$ and
$h_{k+1}=\cdots=h_d=-H$, and take $t>a_k$ so close to $a_k$ that $t<a_{k-1}$ if
$k>2$. Then the $k$-face score is positive, the largest vertex score is 0, and
every face with $2\le j<k$ states scores at most $(j-1)(t-a_j)<0$. A face with
$j>k$ states scores at most $(j-1)(t-a_j)-H(j-k)/j$, and another face with $k$
states at most $(k-1)(t-a_k)-H/k$, so both lose once $H$ is large. The
comparisons are strict, so they hold on an open interval of $t$, and small
changes make the fields strictly ordered without removing the phase.
\end{proof}

\begin{proof}[Proof of Corollary~\ref{cor:weakfield-full-hierarchy}]
Since $\frac1k\sum_{i=1}^k(i-1)^2=(k-1)(2k-1)/6$, the lines
\eqref{eq:weakfield-lines} are $\ell_k(t)=(k-1)t+c_k-H(k-1)(2k-1)/6$, also for
$k=1$, and $\ell_{k+1}(t)-\ell_k(t)=t-t_k$. Extend $c_k$ to the smooth function
$c(x)=(x+\frac12)\log x-\frac{x-1}2\log(4\pi)$ for $x\ge1$. Then
$0<c''(x)=(x-\frac12)/x^2\le\frac12$, and integrating twice over unit intervals
gives $c_{k+2}-2c_{k+1}+c_k\le\frac12$. Hence
\[
 t_{k+1}-t_k=\frac{2H}{3}-(c_{k+2}-2c_{k+1}+c_k)>0.
\]
For $t_{k-1}<t<t_k$, summing the differences $\ell_{j+1}(t)-\ell_j(t)=t-t_j$
shows that line $k$ strictly exceeds every other line. The fields are strictly
decreasing, so the unique best support of size $k$ is $\{1,\ldots,k\}$, and
Theorem~\ref{thm:weakfield-phases} gives the conclusion for finite $N$. The
arithmetic in this proof is checked in Lean.
\end{proof}
